\chapter{Introduction}
This research investigates the application of on-device machine learning to summarise clinical consultations. As a proof-of-concept, the primary objective is to evaluate the feasibility of an offline machine learning pipeline running on Intel AI PC hardware. This aims to reduce an NHS GP's workload by autonomously recording and transcribing consultations, generating clinical summaries, and retrieving medical guidance within an entirely offline environment. 

While existing cloud-based systems such as TORTUS AI, Heidi Health, and Nuance DAX Copilot highlight the utility of automated clinical documentation within healthcare, they rely on networked deployment, extensive integration, and external data processing \cite{tortus_product_2026,heidi_product_2026,dictationdirect_dax_features_2026}. In the context of the NHS, these dependencies introduce significant data governance and infrastructural challenges. Consequently, developing an offline proof-of-concept provides a mechanism to evaluate an alternative deployment architecture designed to address these specific operational constraints.

This project was developed as part of the UCL IXN in partnership with Intel with consultation from Dr. Atia Rafiq, an NHS GP and GP trainer and Dr. Tazeen Ahmed an NHS Rheumatology Care Group Lead.

\section{Problem Overview}
In a consultation, a doctor needs to listen, ask follow-up questions, reason about symptoms, and still produce a record that is accurate for later clinical use. The usual consultation flow is as follows: type and take notes during the appointment, ask follow-up questions, and reconstruct the encounter afterwards. Each of these stages has their drawbacks. Typing during the consultation interrupts eye contact and changes the rhythm of the conversation and writing notes afterwards is slow and risks losing detail.

The main idea of this project is to shift the administrative work into an on-device pipeline. The application records the conversation, transcribes and diarises (labelling sentences as doctor or patient) in real time, produces a clinical summary after the consultation, and searches uploaded guidance documents for relevant medical guidelines. In practice, the output of the consultation contains a transcript, a summary, and a link back to guidance that may help the clinician check what was discussed. This provides a practical basis for evaluating whether an offline workflow can support useful clinical documentation without relying on cloud services.

Running the full pipeline locally introduces its own problems. Large language models are expensive to run. All models must be fast to load and infer enabled by an adequately designed concurrent pipeline ensuring the application remains responsive. Retrieval and summarisation has to be fast, otherwise it becomes one more waiting step in an already busy workflow. The system also has to fit into the constraints of the target platform, a native Windows Store application built with the MSIX packaged project platform and deployed as an MSIX package on Intel AI PC hardware.

\section{Aims}
The project has three main aims:

\begin{itemize}
    \item Evaluate whether the selected models and libraries are fast enough and accurate enough to support an offline workflow in practice.
    \item Build a working proof of concept for offline clinical summarisation on Intel AIPC hardware, with a native Windows user interface and an entirely local inference pipeline.
    \item Produce a codebase and architecture that can be extended after the project, rather than a one-off demonstration that is difficult to maintain.
\end{itemize}

A personal secondary aim is established based on the project requirements. The majority of my development experience throughout my university journey is with Python and VSCode. The technical constraints of this project provide the opportunity to learn the C++ programming language, the Visual Studio environment and the MSIX Packaged Project platform. 

\section{Goals}
From those aims, I defined the following concrete project goals:

\begin{itemize}
    \item Compare candidate pre-trained medical large language models (LLMs) under different quantisation levels and select a summarisation model that is feasible on the target hardware.
    \item Build an Ambient Voice Technology (AVT) pipeline to record and transcribe conversations
    \item Build a retrieval pipeline that allows clinicians to upload medical guidance PDFs and locate relevant sections from the generated consultation summary.
    \item Provide a Windows-native interface for recording, reviewing outputs, selecting and saving external files, and browsing retrieved guidance.
    \item Structure the application so that future work, such as appraisal support or EPR integration, can be added without major rewrites.
\end{itemize}

These goals are intentionally narrower than the long-term ambitions of the product idea. For example, direct EPR integration, multilingual support, and patient-facing summaries are useful extensions, but they are not required to establish whether the core offline pipeline is viable for this proof-of-concept project.

\section{Development Methodology}
This project was set out with two phases, feasibility followed by iterative development on a 1 week cycle. Initially a feasibility study was conducted to determine whether latency and accuracy within the summarisation pipeline could remain within the requirement specification. This involved testing candidate models by measuring inference times, and evaluating performance on a summarisation benchmark.

Development began once technical this feasibility study concluded. Weekly meetings with my external supervisor and regular discussion with my internal supervisor and the NHS clients expanded the small starting feature pool. The original requirements were deliberately narrow, focusing only on local summarisation. Through demonstrations and feedback collection, the scope broadened to include AVT, medical guidance retrieval, and appraisal entries. 

The project is maintained on GitHub, which makes it practical to isolate experiments, preserve stable versions for demonstrations, and revert changes when a promising approach had flaws that required a reset to fix. This mattered most when integrating several highly linked libraries simultaneously with very different potential points of failure.

This flexible development lifecycle suited a proof-of-concept project with several unknowns and regularly evolving requirements. By prioritising high-risk ambiguities in the requirements at the start of development, a stable baseline could be set for demonstrations, thus pre-empting repository refactors due to sudden shifts in requirements.

\section{Overview}
The remainder of this report is organised as follows.

Chapter 2 reviews the background to the project, surveys comparable systems, and explains the choice of models, frameworks, and supporting libraries.

Chapter 3 defines the project requirements in detail and analyses the consequences of those requirements for the system architecture, data model, workflows, and stakeholder responsibilities.

Chapter 4 describes the design and implementation of the application, including the structure of the main components and the reasoning behind the architectural decisions.

Chapter 5 presents the testing strategy, summarises the evidence gathered during development, and identifies the main limitations of the current proof of concept.

Chapter 6 reviews the project against its original goals, critically evaluates the resulting system, and outlines the most important directions for future work.